%%%%%%%%%%%%%%%%%%%%%%%%%%%%%%%%%%%%%%%%%%%%%%%%%%%%%%%%%%%%%%%%%%%%%
%%%%%%%%%%%%%%%%%%%%%%%%%%%%%%%%%%%%%%%%%%%%%%%%%%%%%%%%%%%%%%%%%%%%%
%%%%%%%%%%%%%%%%%%%%%%%%%%%%%%%%%%%%%%%%%%%%%%%%%%%%%%%%%%%%%%%%%%%%%
\section{Bigraded Ginzburg algebras and an invariant polynomial}
\label{sec:euler}

The object of this section is to develop the study of the graded Euler pairing of a differential bigraded Ginzburg algebra associated to a graded quiver with potential. Two conceptually important facts we establish are \Cref{thm:intrinsic-euler}, showing that the determinant of the graded Euler pairing is a well-defined invariant of the dilation class, and \Cref{prop:mutation-euler}, showing that it is invariant under graded QP mutations. Several examples are provided in \Cref{ssec:two_examples} and \Cref{ssec:tree_quivers}.

%%%%%%%%%%%%%%%%%%%%%%%%%%%%%%%%%%%%%%%%%%%%%%%%%%%%%%%%%%%%%%%%%%%%% 
\subsection{The internal grading of the Ginzburg dg algebra of a graded QP}

Let \((Q,W)\) be a QP and $\Gamma(Q,W)$ its complete 3-Calabi-Yau Ginzburg dg algebra, cf.~\cite[Section 4.2]{Ginzburg} and \cite[Section~2.6]{KellerYang}. That is, $\Gamma(Q,W)=\widehat{\kk{\tilde{Q}}}$ is the completed path algebra of the doubled quiver $\tilde{Q}$, where we take the completion as in \cite{KellerYang}. The original arrows $a:i\to j$, $a\in Q_1$, are graded in cohomological degree 0, the dual arrows $a^*:j\to i$, $a^*\in \tilde{Q}_1$, are given cohomological degree -1 and the loop arrows $t_i:i\to i$, dual to the vertices, are given cohomological degree -2. The differential $\Gamma(Q,W)$ is declared on these generators to be

\begin{align*}
 \partial(a)=0,\quad \partial(a^*)=\partial_aW,\quad \partial(t_i)=e_i\sum_{a\in Q_1}[a,a^*]e_i.
\end{align*}

If $(Q,W,d)$ is a graded QP, then the dg algebra $\Gamma(Q,W)$ can be further endowed with an internal grading, a.k.a.~an Adams grading, as follows:

\begin{definition}[Bigraded Ginzburg dga of a graded QP]\label{def:bigradedGinzburg} Let \((Q,W,d)\) be a graded QP. By definition, the internal (or Adams) grading on the Ginzburg dg algebra $\Gamma(Q,W)$ of $(Q,W)$ is given by:
\begin{enumerate}
    \item The original arrows $a:i\to j$ in $Q_1\sse \tilde{Q}_1$ are assigned internal degree $d(a)$,
    \item The dual arrows $a^*:j\to i$ in $\tilde{Q}_1$ are assigned internal degree $1-d(a)$,
    \item The loop arrows $t_i:i\to i$ in $\tilde{Q}_1$ are assigned internal degree $1$.
\end{enumerate}
The resulting differential bigraded Ginzburg algebra is denoted by $\Gamma(Q,W,d)$. In table form, the bidegrees of the generators of $\Gamma(Q,W,d)$ are
\begin{center}
\begin{tabular}{c|c|c}\label{table:gradings}
 generator & cohomological degree & internal degree\\
\midrule
 \(a:i\to j\) & \(0\) & \(d(a)\)\\
 \(a^*:j\to i\) & \(-1\) & \(1-d(a)\)\\
 \(t_i:i\to i\) & \(-2\) & \(1\)
\end{tabular}
\end{center}
\hfill$\Box$
\end{definition}

\noindent Since \(W\) has internal degree one, the differential in $\Gamma(Q,W,d)$ has bidegree
\((1,0)\). Note also that sending all arrows to zero produces an augmentation $\Gamma(Q,W,d)\to R_Q$.

%%%%%%%%%%%%%%%%%%%%%%%%%%%%%%%%%%%%%%%%%%%%%%%%%%%%%%%%%%%%%%%%%%%%% 
\subsection{Module categories of the bigraded Ginzburg algebra}\label{ssec:module_categories_bigraded} The topology of $\Gamma(Q,W)$ -- and in our case $\Gamma(Q,W,d)$ -- should be considered, as in \cite[Appendix]{KellerYang} and see also \cite{Keller2026_GradedPreprint}. The complete Ginzburg algebra $\Gamma(Q,W)$ is pseudocompact for the $\mathfrak{m}_{\tilde{Q}}$-adic topology, and we consider {\it continuous} right dg modules over it, cf.~\cite[Appendix]{KellerYang}. In the graded case, for $\Gamma:=\Gamma(Q,W,d)$, the internal grading is incorporated by considering continuous dg modules $M$ over $\Gamma(Q,W)$ equipped with closed subcomplexes $M^{(r)}$, $r\in\Z$, and a topological product decomposition
\[
  M\cong\prod_{r\in\Z}M^{(r)},\quad\mbox{such that}\quad
  M^{(r)}\Gamma^{(s)}\subset M^{(r+s)}\quad\mbox{and}\quad
  \partial_M(M^{(r)})\subset M^{(r)}.
\]
In addition, we require a neighborhood basis by homogeneous dg-submodules of finite total codimension.  Equivalently, such a continuous differential bigraded module $M$ over $\Gamma(Q,W,d)$ is an inverse limit of finite-dimensional \emph{internally graded} discrete dg modules and degree-preserving transition maps. As before, the cohomological differential therefore has bidegree $(1,0)$, and a morphism in the graded category is continuous of bidegree $(0,0)$.\\

Localizing the homotopy category of these continuous internally graded dg modules over $\Gamma(Q,W,d)$ at quasi-isomorphisms gives the derived category $\D^{\gr}_{\cont}(\Gamma)$. For $M,N\in D^{\mathrm{gr}}_{\mathrm{cont}}(\Gamma)$, we define
\[
\operatorname{Ext}^p_\Gamma(M,N)_r
:=
H^p\!\left(
\operatorname{RHom}_\Gamma(M,N)_{-r}
\right)
\cong
\operatorname{Hom}_{D^{\mathrm{gr}}_{\mathrm{cont}}(\Gamma)}
\bigl(M,N[p]\langle-r\rangle\bigr).
\]
Here $[p]$ denotes the cohomological shift and $\langle r\rangle$
the internal shift. Our internal-shift convention is $(N\langle r\rangle)_p:=N_{p+r}$. Thus an element of internal degree $p$ in $N$ has internal degree
$p-r$ in $N\langle r\rangle$. In particular,
$\Bbbk\langle r\rangle$ is concentrated in ordinary internal degree
$-r$. The subscript $r$ on
$\operatorname{Ext}^p_\Gamma(M,N)_r$ records the shift of the internal
weight, i.e.~ the negative of the (standard) internal
degree of the corresponding homogeneous morphism.\\

We will be working with the following subcategories of $\D^{\gr}_{\cont}(\Gamma)$:

\begin{definition}\label{def:categories_over_bigraded}
Let $(Q,W,d)$ be a graded QP and $\D^{\gr}_{\cont}(\Gamma)$ its derived category of continuous internally graded dg modules over $\Gamma:=\Gamma(Q,W,d)$. By definition:

\begin{enumerate}
    \item The perfect category $\per^{\gr}(\Gamma)$ is the smallest idempotent-complete triangulated subcategory of $\D^{\gr}_{\cont}(\Gamma)$ containing the free module and all of its internal shifts:
\begin{equation}\label{eq:perfect-category}
  \per^{\gr}(\Gamma)
  :=\thick\set{\Gamma\angles r\suchthat r\in\Z}
  =\thick\set{e_i\Gamma[p]\angles r
      \suchthat i\in Q_0,\ p,r\in\Z}.
\end{equation}

\item The finite-dimensional category $\D^{\gr}_{\fd}(\Gamma)$ is
\[
  \D^{\gr}_{\fd}(\Gamma)
  :=\set{M\in\D^{\gr}_{\cont}(\Gamma)\suchthat
  \sum_p\dim_\kk H^p(M)<\infty}.
\]
\end{enumerate}
\end{definition}

For $\per^{\gr}(\Gamma)$, its objects are precisely those objects which are isomorphic in the derived category to direct summands of finite-cell, topologically semi-free dg modules.\footnote{That is, we start with finitely many shifted modules $e_i\Gamma[p]\angles r$, form finitely many cones, and then take finite direct sums and direct summands.} For $\D^{\gr}_{\fd}(\Gamma)$ note that, since the total cohomologies in $\D^{\gr}_{\fd}(\Gamma)$ are finite-dimensional, only finitely many internal weights can occur in its objects.\\

Let us now focus on certain objects of $\D^{\gr}_{\cont}(\Gamma)$, namely, the simple modules associated to the vertices of $Q$. By definition, the simple right dg module $S_i$ over $\Gamma$ associated to the vertex $i\in Q_0$ is
\begin{equation}\label{eq:Si}
  S_i:=e_iR_Q,
\end{equation}

\noindent where $\Gamma$ acts via the augmentation $\Gamma\lr R_Q$ that sends all arrows to zero. See \Cref{ssec:two_examples} for some examples of such simple modules. Such simple modules $S_i$ are concentrated in cohomological and internal degree zero.

\begin{proposition}[Finiteness of Ext groups between simples]\label{prop:ext-finite} Let $(Q,W,d)$ be a graded QP and $\Gamma=\Gamma(Q,W,d)$ its associated Ginzburg algebra. Then:
\begin{enumerate}
    \item Each simple $S_i$ belongs to $\D^{\gr}_{\fd}(\Gamma)\cap\per^{\gr}(\Gamma)$.
    
    \item For every pair of vertices $i,j\in Q_0$, the vector spaces $\Ext^{p}_{\Gamma}(S_i,S_j)$ vanish for $p\notin\{0,1,2,3\}$. In particular, $\Ext^{p}_{\Gamma}(S_i,S_j)_r$ vanish for $p\notin\{0,1,2,3\}$

    \item For every pair of vertices $i,j\in Q_0$, the vector spaces $\Ext^{p}_{\Gamma}(S_i,S_j)$ are finite-dimensional for all $p\in\Z$. In particular, so are $\Ext^{p}_{\Gamma}(S_i,S_j)_r$.
\end{enumerate}
\end{proposition}

\begin{proof}
For Part (1), note that the module $S_i$ is one-dimensional by \Cref{eq:Si}, and hence it necessarily lies in $\D^{\gr}_{\fd}(\Gamma)$. For its perfectness, consider the explicit cofibrant resolution of $S_i$ constructed in \cite[Section~2.14]{KellerYang} and note that their construction is homogeneous for the additional internal grading because the Ginzburg differential has bidegree $(1,0)$. (See also \Cref{ssec:two_examples} below for some explicit instances of these cofibrant resolutions.) Since the resolution is a finite semi-free complex built from the projectives $e_j\Gamma$, it follows that $S_i$ is perfect and thus Part (1) is established.\\

\noindent For Part (2), it follows from \cite[Lemma~2.15]{KellerYang} that a basis of $\Hom_{\D(\Gamma)}(S_i,S_j[p])$ is given by generators of the Ginzburg quiver from $j$ to $i$ of cohomological degree $-p+1$, together with the identity when $i=j$ and $p=0$.  Since the Ginzburg quiver has arrows only in cohomological degrees $0,-1,-2$, Part (2) follows. Finally, finiteness of $Q$ gives finite-dimensionality and finite internal-weight support, concluding Part (3).
\end{proof}

%%%%%%%%%%%%%%%%%%%%%%%%%%%%%%%%%%%%%%%%%%%%%%%%%%%%%%%%%%%%%%%%%%%%% 
\subsection{Graded Euler pairing}\label{ssec:graded_Euler_pairing} Let $(Q,W,d)$ be a graded QP and $\Gamma=\Gamma(Q,W,d)$ its associated differential bigraded Ginzburg algebra, as in \Cref{def:bigradedGinzburg}. If the internal grading is not taken into account, the Euler form in the Grothendieck group of the derived category of finite-dimensional modules over $\Gamma(Q,W)$ is an integer, to be interpreted as an element of the Grothendieck group of the trivial group. By incorporating the internal grading of $\Gamma$, as in \Cref{def:bigradedGinzburg}, the Euler form in the Grothendieck group of $\D^{\gr}_{\fd}(\Gamma)$ should be valued in $K_0(\mathbb{G}_m)\cong\Z[t,t^{-1}]$, i.e.~it should be a $\mathbb{G}_m$-equivariant (or $\Z$-graded) Euler characteristic.\footnote{It is possible to phrase the study of $\Z$-graded QPs in terms of $\mathbb{G}_m$ actions $\rho:\mathbb{G}_m\lr\mbox{Aut}_{R_Q}(\widehat{\kk Q})$ on the completed path algebra which act via $\rho_z(W)=zW$. In that viewpoint, it is natural to consider an equivariant Euler characteristic.} Here we use the coordinate $t$ to correspond with the internal shift, so that $[M\langle r\rangle]=t^r[M]$.\\

The simple modules $S_i$ are such that the Grothendieck group $K_0(\D^{\gr}_{\fd}(\Gamma))$ is a free $\Z[t,t^{-1}]$-module with basis $[S_i]$. Thus we can express the graded Euler pairing in terms of this basis of simple classes:

\begin{definition}[Graded Euler pairing in the simple basis]\label{def:euler}
Let $(Q,W,d)$ be a graded QP and $\Gamma=\Gamma(Q,W,d)$ its associated differential bigraded Ginzburg algebra. By definition, the graded Euler matrix of
\((Q,W,d)\) is
\begin{equation}\label{eq:euler-ext}
 \Eul_{(Q,W,d)}(t)_{ij}
 :=\sum_{p,r\in\Z}(-1)^p
 \dim_\kk\Ext^{p}_{\Gamma}(S_i,S_j)_r\cdot t^{r}.
\end{equation}
\hfill$\Box$\end{definition}
\noindent That is, the entry $\Eul_{(Q,W,d)}(t)_{ij}$ in \Cref{def:euler} is the graded Euler characteristic of $\RHom_{\Gamma}(S_i,S_j)$. It is written in
the internal-shift coordinate $t$, so that a copy of
$\Bbbk\langle r\rangle$ contributes $t^r$. The expansion \eqref{eq:euler-ext} can also be written as

\begin{equation}\label{eq:euler-ext2}
 \Eul_{(Q,W,d)}(t)_{ij}
 =\sum_{p\in\Z}(-1)^p
 \mbox{ch}_t\Ext^{p}_{\Gamma}(S_i,S_j)=\sum_{p=0}^3(-1)^p
 \mbox{ch}_t\Ext^{p}_{\Gamma}(S_i,S_j)
\end{equation}
where the second equality in \eqref{eq:euler-ext2} is \Cref{prop:ext-finite}.(2). In \eqref{eq:euler-ext2}, $ \mbox{ch}_t\Ext^{p}_{\Gamma}(S_i,S_j)$ denotes the equivariant character for the corresponding $\Gm$-action on $\Ext^{p}_{\Gamma}(S_i,S_j)$. Specifically, we used the equivariant character
\[
\operatorname{ch}_t(V)
:=
\sum_{r\in\mathbb Z}
\dim_\Bbbk\!\left(V_{-r}\right)t^r
\]
for a graded vector space $V$. Equivalently, if $V\cong\bigoplus_{r\in\mathbb Z}V_r\otimes_\Bbbk\Bbbk\langle r\rangle$ then $\operatorname{ch}_t(V)=\sum_r(\dim_\Bbbk V_r)t^r$. Such graded (or equivariant) Euler characteristic \eqref{eq:euler-ext} is additive in distinguished triangles and therefore defines an Euler pairing on the equivariant Grothendieck group. It is a sesquilinear form for the involution $t\to t^{-1}$. Note that \Cref{prop:ext-finite} implies that the entries \eqref{eq:euler-ext} are Laurent polynomials in the $t$-variable, as there are only finitely many non-vanishing coefficients and those are themselves finite. To ease notation we often denote $\Eul_{ij}(t):=\Eul_{(Q,W,d)}(t)_{ij}$ if $(Q,W,d)$ is clear by context.

\begin{remark} (1) The more ambitious prospect of considering the bivariate polynomial
\begin{align*}
  \Poin_{\Gamma}(u,t)_{ij}
  &=\sum_{p,r\in\Z}
  \dim_\kk\Ext^{p}_{\Gamma}(S_i,S_j)_ru^pt^{r}
\end{align*}
does not work for our purposes, as its behavior under equivalences and mutations is not by congruence. In part, this failure of functoriality for the above bivariate polynomial comes from the fact that such double-Euler characteristic is not additive in distinguished triangles.\\

\noindent (2) Some formulas in the quiver literature index the Euler matrix by arrows
\(i\to j\) rather than by \(\Ext(S_i,S_j)\): the resulting matrix is then \(\Eul(t)^{\mathsf T}\). Note that, nevertheless, their determinants are unaffected.
\hfill$\Box$\end{remark}

It will be useful to have a more hands-on description of the graded Euler pairing directly in terms of the quiver and its grading. We establish combinatorial formulas for the entries as follows:

\begin{proposition}[Combinatorial Euler pairing in the simple basis]\label{prop:euler-formula}
Let $(Q,W,d)$ be a graded QP. In the notation above,
\begin{equation}\label{eq:euler-formula}
 \boxed{
 \Eul_{ij}(t)
 =(1-t)\delta_{ij}
 -\sum_{a:j\to i}t^{d(a)}
 +\sum_{a:i\to j}t^{1-d(a)}.}
\end{equation}
\end{proposition}

\begin{proof}
It suffices to use \cite[Lemma~2.15]{KellerYang}, which computes
\(\Hom(S_i,S_j[p])\) from the generators of the Ginzburg quiver that are directed from
\(j\) to \(i\). The contributions are thus as follows:

\begin{enumerate}
    \item The identity contributes \(+1\) for \(i=j,p=0\), corresponding to the summand $\delta_{ij}$ in \Cref{eq:euler-formula}.\\
    \item An original arrow \(a:j\to i\), $a\in Q_1$, has cohomological degree zero and internal degree \(d(a)\). The corresponding element of
$\operatorname{Ext}^1_\Gamma(S_i,S_j)$ has (standard) internal degree
$-d(a)$ and therefore contributes $-t^{d(a)}$. These are the contributions to the first (negative) sum in \Cref{eq:euler-formula}.\\
    \item If \(a:i\to j\), the opposite arrow \(a^*:j\to i\) has cohomological degree \(-1\) and internal degree \(1-d(a)\). The corresponding element of $\operatorname{Ext}^2_\Gamma(S_i,S_j)$ has standard internal degree $d(a)-1$ and hence contributes \(+t^{1-d(a)}\). These are the contributions to the second sum in \Cref{eq:euler-formula}.\\
    \item Finally, a loop \(t_i:i\to i\) has internal degree 1 and the
corresponding element of
$\operatorname{Ext}^3_\Gamma(S_i,S_i)$ has internal degree
$-1$. Thus a loop contributes $-t$. This is the term $-t\delta_{ij}$ in \Cref{eq:euler-formula}.
\end{enumerate}
Thus adding these four contributions gives \eqref{eq:euler-formula}.
\end{proof}

\noindent Note that \Cref{eq:euler-formula} depends only on the graded arrow bimodule and not
on the coefficients of \(W\). It is particularly apparent from \Cref{eq:euler-formula} that the ungraded Euler pairing, obtained by setting $t=1$, yields the exchange matrix of the quiver $Q$. In this sense, the graded Euler pairing can be understood as a deformation of the exchange matrix with parameter $t$.\\

The key point now is to understand how the equivariant Euler characteristics in \Cref{eq:euler-ext} behave in the following situations:\\

\begin{itemize}
    \item[(a)] The first task is to show that, while $\Eul_{(Q,W,d)}(t)_{ij}$ is defined using $(Q,W,d)$, the conjugacy class of the matrix $(\Eul_{(Q,W,d)}(t)_{ij})\in\mbox{Mat}(\Z[t,t^{-1}])$ is actually a well-defined invariant associated to any dilation class $\delta\in\Gclass[Q,W]$. In fact, more is true, as its monomial-conjugacy class will be well-defined for a dilation class.\\

    \item[(b)] The second task is to show that the congruence class of the matrix $(\Eul_{(Q,W,d)}(t)_{ij})$ is an invariant of the graded QP-mutation class of $(Q,W,d)$.\\
\end{itemize}

\noindent These verifications can be made using categorical arguments or by explicitly computing with \Cref{eq:euler-formula}. \Cref{sec:euler-dilation} implements task $(a)$ and \Cref{sec:euler-mutation} addresses task $(b)$.

%%%%%%%%%%%%%%%%%%%%%%%%%%%%%%%%%%%%%%%%%%%%%%%%%%%%%%%%%%%%%%%%%%%%% 
\subsection{Two examples}\label{ssec:two_examples} Let us provide the details for how to compute in some examples the graded Euler pairing \eqref{eq:euler-ext}, in the basis of simple classes, as defined in \Cref{def:euler}.\\

First, let $(Q,W)$ be the 4-cycle QP, with quiver 
\begin{center}
\begin{tikzcd}[row sep=large, column sep=large]
    4 \arrow[d, "d"'] & 3 \arrow[l, "c"'] \\
    1 \arrow[r, "a"'] & 2 \arrow[u, "b"']
\end{tikzcd}
\end{center}
and potential $W = dcba$. For the arrow grading $\deg:Q_1\lr\Z$, let us assign internal degrees as
$$\deg(a)=1,\deg(b)=\deg(c)=\deg(d)=0,$$ and thus \Cref{def:bigradedGinzburg} dictates $$\deg(a^*)=0,\deg(b^*)=\deg(c^*)=\deg(d^*)=1,$$
while all the loops $t_i$ have $\deg(t_i)=1$. Note that the potential $W$ is indeed homogeneous of degree 1. This defines the differential bigraded Ginzburg algebra $\Gamma(Q,W,\deg)$ in this case.\\

\noindent Both the proof of \Cref{prop:ext-finite} and \Cref{prop:euler-formula} implicitly use cofibration resolutions of the simple modules $S_i$, as built in \cite[Section 2.14]{KellerYang}. Let us recall here that, as the proof of \cite[Section 3.12]{KellerYang} illustrates, such cofibrant resolution $pS_i$ of $S_i$ is constructed as a shifted mapping cone, with underlying vector space:

\begin{equation}\label{eq:cofibrant_resolution}
pS_i := \Sigma^3 P_i \oplus \left(\bigoplus_{\alpha \in Q_1: s(\alpha)=i} \Sigma^2 P_{t(\alpha)}\right) \oplus \left(\bigoplus_{\beta \in Q_1: t(\beta)=i} \Sigma P_{s(\beta)} \oplus P_i\right),
\end{equation}
where $(P_i,\partial_{P_i})$ denotes the projective dg module associated to the vertex $i\in Q_0$, and $\Sigma$ denotes suspension. The corresponding differential $\partial_{pS_i}$ acting on the vector space \eqref{eq:cofibrant_resolution} is given by the following block lower-triangular matrix, where the entries operate via left multiplication:
\begin{equation}\label{eq:cofibration_resolution_differential}
\partial_{pS_i} := \begin{pmatrix} 
\partial_{\Sigma^3 P_i} & 0 & 0 & 0 \\ 
\alpha & \partial_{\Sigma^2 P_{t(\alpha)}} & 0 & 0 \\ 
-\beta^* & -\partial_{\alpha\beta}W & \partial_{\Sigma P_{s(\beta)}} & 0 \\ 
t_i & \alpha^* & \beta & \partial_{P_i} 
\end{pmatrix}.
\end{equation}
Recall that we use the internal grading convention $(M\langle r \rangle)_{q} = M_{q+r}$,  cf.~\Cref{ssec:module_categories_bigraded}. In particular, with this convention, if a component of a differential is given by left multiplication by a homogeneous element of internal degree $s$, then the corresponding source projective module must be shifted by
$\langle-s\rangle$ in order for that component to have
internal degree zero. The minimal semi-free cofibrant resolutions $pS_i$ for the four simples in our example of $(Q,W)$ then are:
\begin{align*}
pS_1 &= P_1[3]\langle -1 \rangle \oplus P_2[2]\langle 0 \rangle \oplus P_4[1]\langle 0 \rangle \oplus P_1[0]\langle 0 \rangle \\
pS_2 &= P_2[3]\langle -1 \rangle \oplus P_3[2]\langle -1 \rangle \oplus P_1[1]\langle -1 \rangle \oplus P_2[0]\langle 0 \rangle \\
pS_3 &= P_3[3]\langle -1 \rangle \oplus P_4[2]\langle -1 \rangle \oplus P_2[1]\langle 0 \rangle \oplus P_3[0]\langle 0 \rangle \\
pS_4 &= P_4[3]\langle -1 \rangle \oplus P_1[2]\langle -1 \rangle \oplus P_3[1]\langle 0 \rangle \oplus P_4[0]\langle 0 \rangle
\end{align*}


\begin{table}[h]
\centering
\caption{Non-vanishing internally graded Ext groups $\Ext^p(S_i,S_j)$ for the non-degenerate 4-cycle graded QP $(Q,W,\deg)$.}
\label{table:Ext_groups1}
\begin{tabular}{@{}l c c c c@{}}
\toprule
Pair $(S_i, S_j)$ & $p=0$ & $p=1$ & $p=2$ & $p=3$ \\ 
\midrule
$(S_1, S_1)$ & $k\langle 0 \rangle$ & 0 & 0 & $k\langle 1 \rangle$ \\
$(S_2, S_2)$ & $k\langle 0 \rangle$ & 0 & 0 & $k\langle 1 \rangle$ \\
$(S_3, S_3)$ & $k\langle 0 \rangle$ & 0 & 0 & $k\langle 1 \rangle$ \\
$(S_4, S_4)$ & $k\langle 0 \rangle$ & 0 & 0 & $k\langle 1 \rangle$ \\
\addlinespace
$(S_1, S_2)$ & 0 & 0 & $k\langle 0 \rangle$ & 0 \\
$(S_2, S_3)$ & 0 & 0 & $k\langle 1 \rangle$ & 0 \\
$(S_3, S_4)$ & 0 & 0 & $k\langle 1 \rangle$ & 0 \\
$(S_4, S_1)$ & 0 & 0 & $k\langle 1 \rangle$ & 0 \\
\addlinespace
$(S_2, S_1)$ & 0 & $k\langle 1 \rangle$ & 0 & 0 \\
$(S_3, S_2)$ & 0 & $k\langle 0 \rangle$ & 0 & 0 \\
$(S_4, S_3)$ & 0 & $k\langle 0 \rangle$ & 0 & 0 \\
$(S_1, S_4)$ & 0 & $k\langle 0 \rangle$ & 0 & 0 \\
\bottomrule
\end{tabular}
\end{table}

\noindent By considering the homology of the morphism complex $\mathcal{H}om_{\Gamma}(pS_i, S_j)$, a summand $P_j[h]\langle r\rangle$ contributes a one-dimensional
vector space in internal degree $r$, equivalently to the component $\operatorname{Ext}^p_\Gamma(S_i,S_j)_{-r}$. The resulting non-zero graded Ext groups between the simples of $\Gamma(Q,W,d)$ are computed in \Cref{table:Ext_groups1}.\\

In order to compute the graded Euler pairing, as in \Cref{def:euler}, let $t$ be the formal parameter keeping track of the shift in internal degree. From \Cref{table:Ext_groups1}, the equivariant Euler pairing in the basis of the simples gives the matrix:

\[
\Eul_{(Q,W,d)}(t) = \begin{pmatrix} 
1 - t & 1 & 0 & -1 \\ 
-t & 1 - t & t & 0 \\ 
0 & -1 & 1 - t & t \\ 
t & 0 & -1 & 1 - t 
\end{pmatrix}.
\]
This concludes this first example. For our second example, we will consider the following QP $(Q',W')$, which adds two vertices to $(Q,W)$:

\begin{center}
$Q':=$\begin{tikzcd}[row sep=large, column sep=large]
    4 \arrow[d, "d"'] & 3 \arrow[l, "c"'] \arrow[r, "g"] & 6 \arrow[d, "f"] \\
    1 \arrow[r, "a"'] & 2 \arrow[u, "b"'] & 5 \arrow[l, "e"]
\end{tikzcd}
\end{center}
with potential $W' := dcba + efgb$. The underlying quiver is of finite mutation type $E_6$. The internal grading $\deg'$ for the graded QP $(Q',W',\deg')$ is that of $(Q,W,d)$ now extended such that
$$\deg(e)=1,\quad \deg(f)=\deg(g)=0,$$
and so their duals have degrees
$$\deg(e^*)=0,\quad \deg(f^*)=\deg(g^*)=1.$$ 

\noindent Following \eqref{eq:cofibrant_resolution} as above, the cofibrant resolutions $pS_i$ for the simples of $Q'$ have underlying vector spaces:
\begin{align*}
pS_1 &= P_1[3]\langle -1 \rangle \oplus P_2[2]\langle 0 \rangle \oplus P_4[1]\langle 0 \rangle \oplus P_1[0]\langle 0 \rangle \\
pS_2 &= P_2[3]\langle -1 \rangle \oplus P_3[2]\langle -1 \rangle \oplus P_1[1]\langle -1 \rangle \oplus P_5[1]\langle -1 \rangle \oplus P_2[0]\langle 0 \rangle \\
pS_3 &= P_3[3]\langle -1 \rangle \oplus P_4[2]\langle -1 \rangle \oplus P_6[2]\langle -1 \rangle \oplus P_2[1]\langle 0 \rangle \oplus P_3[0]\langle 0 \rangle \\
pS_4 &= P_4[3]\langle -1 \rangle \oplus P_1[2]\langle -1 \rangle \oplus P_3[1]\langle 0 \rangle \oplus P_4[0]\langle 0 \rangle \\
pS_5 &= P_5[3]\langle -1 \rangle \oplus P_2[2]\langle 0 \rangle \oplus P_6[1]\langle 0 \rangle \oplus P_5[0]\langle 0 \rangle \\
pS_6 &= P_6[3]\langle -1 \rangle \oplus P_5[2]\langle -1 \rangle \oplus P_3[1]\langle 0 \rangle \oplus P_6[0]\langle 0 \rangle
\end{align*}

\begin{table}[h]
\centering
\caption{Non-vanishing internally graded Ext groups $\Ext^p(S_i,S_j)$ for our second example, the non-degenerate graded QP $(Q',W',\deg')$.}
\label{table:Ext_groups2}
\begin{tabular}{@{}l c c c c@{}}
\toprule
Pair $(S_i, S_j)$ & $p=0$ & $p=1$ & $p=2$ & $p=3$ \\ 
\midrule
All $(S_i, S_i)$ & $k\langle 0 \rangle$ & 0 & 0 & $k\langle 1 \rangle$ \\
\addlinespace
$(S_1, S_2)$ & 0 & 0 & $k\langle 0 \rangle$ & 0 \\
$(S_2, S_3)$ & 0 & 0 & $k\langle 1 \rangle$ & 0 \\
$(S_3, S_4)$ & 0 & 0 & $k\langle 1 \rangle$ & 0 \\
$(S_4, S_1)$ & 0 & 0 & $k\langle 1 \rangle$ & 0 \\
$(S_5, S_2)$ & 0 & 0 & $k\langle 0 \rangle$ & 0 \\
$(S_3, S_6)$ & 0 & 0 & $k\langle 1 \rangle$ & 0 \\
$(S_6, S_5)$ & 0 & 0 & $k\langle 1 \rangle$ & 0 \\
\addlinespace
$(S_2, S_1)$ & 0 & $k\langle 1 \rangle$ & 0 & 0 \\
$(S_3, S_2)$ & 0 & $k\langle 0 \rangle$ & 0 & 0 \\
$(S_4, S_3)$ & 0 & $k\langle 0 \rangle$ & 0 & 0 \\
$(S_1, S_4)$ & 0 & $k\langle 0 \rangle$ & 0 & 0 \\
$(S_2, S_5)$ & 0 & $k\langle 1 \rangle$ & 0 & 0 \\
$(S_6, S_3)$ & 0 & $k\langle 0 \rangle$ & 0 & 0 \\
$(S_5, S_6)$ & 0 & $k\langle 0 \rangle$ & 0 & 0 \\
\bottomrule
\end{tabular}
\end{table}

The corresponding differentials, as above, follow from \eqref{eq:cofibration_resolution_differential}. The non-vanishing graded Ext groups between the simples of $\Gamma(Q',W',d')$ are summarized in \Cref{table:Ext_groups2}. Consequently, the graded Euler pairing for $(Q',W',\deg')$ in the basis of simple classes is given by:

\[
\Eul_{(Q',W',\deg')}(t) = \begin{pmatrix} 
1 - t & 1 & 0 & -1 & 0 & 0 \\ 
-t & 1 - t & t & 0 & -t & 0 \\ 
0 & -1 & 1 - t & t & 0 & t \\ 
t & 0 & -1 & 1 - t & 0 & 0 \\
0 & 1 & 0 & 0 & 1 - t & -1 \\
0 & 0 & -1 & 0 & t & 1 - t
\end{pmatrix}
\]

\noindent The determinant of $\Eul_{(Q',W',\deg')}(t)$ exactly recovers the Alexander polynomial of the $(3,4)$-torus knot:
\begin{equation}\label{eq:E6_determinant_Alexander}
\det(\Eul_{(Q',W',\deg')}(t)) = t^{6} - t^{5} + t^{3} - t + 1 = \Delta_{3,4}(t).
\end{equation}

\noindent To motivate such connection, note that the $(3,4)$-torus knot is the link of the $E_6$ simple plane curve singularity. By \cite[Theorem 18.3]{FPST}, the quiver associated to any real Morsification of such $E_6$ singularity is of mutation type $E_6$. The example of the $E_6$ simple plane curve singularity and its quivers will be further discussed in \Cref{ssec:examples}. The equality in \eqref{eq:E6_determinant_Alexander} is an instance of the more general phenomenon proven in \Cref{cor:cokernel_module_is_Alexander_algebraic_links}.

%%%%%%%%%%%%%%%%%%%%%%%%%%%%%%%%%%%%%%%%%%%%%%%%%%%%%%%%%%%%%%%%%%%%% 
\subsection{Graded Euler pairing and dilation classes}\label{sec:euler-dilation} Let $(Q,W,d)$ be a graded QP and $(\Eul_{ij}(t))$ the matrix of graded Euler characteristics in the simple basis, as in \Cref{def:euler}. By \Cref{def:G1}, a dilation is an equivalence class under vertex gauging, homogeneous stabilizations and graded right-equivalences. Let us study how $(\Eul_{ij}(t))$ behaves under these equivalence relations:

\begin{lemma}[Effect of vertex gauge]\label{prop:euler-gauge}
Let $(Q,W,d)$ be a graded QP, $h:Q_0\lr\Z$ a vertex grading, and \(d^h\) the result of applying a vertex-gauge to $d$, as in \eqref{eq:gauge}. Consider the diagonal matrices
\[
 D_h(t)=\diag(t^{h(i)})_{i\in Q_0}.
\]
Then the resulting matrix $\Eul_{(Q,W,d^h)}(t)$ of graded Euler characteristics is given by
\begin{equation}\label{eq:euler-gauge}
 \Eul_{(Q,W,d^h)}(t)
 =D_h(t)\Eul_{(Q,W,d)}(t)D_h(t)^{-1}.
\end{equation}
\end{lemma}

\begin{proof}
This follows from \Cref{prop:euler-formula}. Indeed, for an arrow \(a:j\to i\),
\[
 t^{d^h(a)}=t^{h(i)-h(j)}t^{d(a)}.
\]
For an arrow \(a:i\to j\),
\[
 t^{1-d^h(a)}=t^{h(i)-h(j)}t^{1-d(a)}.
\]
Thus every \((i,j)\)-entry of \Cref{eq:euler-formula} is multiplied by
\(t^{h(i)-h(j)}\), which implies
\Cref{eq:euler-gauge}.
\end{proof}

\begin{lemma}[Effect of homogeneous stabilization]\label{lem:trivial-cancel}
Let $(Q,W,d)$ be a graded QP. Then, adjoining or deleting a homogeneous trivial QP summand to $(Q,W,d)$ does not change the Euler matrix, i.e.~ the resulting Euler matrix is equal to $\Eul_{(Q,W,d)}(t)$.
\end{lemma}

\begin{proof}
It suffices to consider arrows
\[
 x:i\to j,
 \qquad
 y:j\to i,
 \qquad\mbox{with }
 d(x)+d(y)=1,
\]
and the potential \(yx\).  In the \((j,i)\)-entry, \(x\) contributes
\(-t^{d(x)}\) and \(y\) contributes
\(+t^{1-d(y)}=+t^{d(x)}\), and so these contributions cancel. The two contributions in the
\((i,j)\)-entry cancel similarly, and no other entries change, thus proving that $\Eul_{(Q,W,d)}(t)$ remains the same.
\end{proof}

\begin{lemma}[Effect of graded right-equivalence]\label{lem:right-equivalence}
Let $(Q,W,d)$ be a graded QP. Then, applying a graded right-equivalence does not change the Euler matrix, i.e.~ the resulting Euler matrix is equal to $\Eul_{(Q,W,d)}(t)$.
\end{lemma}

\begin{proof}
Indeed, under a graded right-equivalence, the induced map on
\(\m/\m^2\) is an isomorphism of graded vertex-bimodules.  Therefore, for
every ordered pair \((i,j)\), it preserves the multiset of arrow degrees
from \(i\) to \(j\). Thus Formula \eqref{eq:euler-formula} gives literal equality
of the matrices.
\end{proof}

Lemmas \ref{prop:euler-gauge}, \ref{lem:trivial-cancel} and \ref{lem:right-equivalence} together imply:

\begin{corollary}[Graded Euler pairing and dilations]\label{thm:intrinsic-euler}
Let $(Q,W)$ be a QP and \(\delta\in\Gclass[Q,W]\) a dilation class. Then the following holds:
\begin{enumerate}
    \item The conjugation class of
\(\Eul_{(Q',W',d')}(t)\) is independent of the representative $(Q',W',d')$ of the dilation class $\delta\in\Gclass[Q,W]$. That is, if $(Q_1,W_1,d_1),(Q_2,W_2,d_2)$ represent the same class in $\Gclass[Q,W]$ then $\Eul_{(Q_1,W_1,d_1)}(t)$ and $\Eul_{(Q_2,W_2,d_2)}(t)$ are conjugate by a matrix in $\GL_{|Q_0|}(\Z[t,t^{-1}])$.\\

\item The isomorphism type of the cokernel $\coker\Eul_\delta(t)$, considered as a $\Z[t,t^{-1}]$-module, is a well-defined invariant of the dilation class $\delta\in\Gclass[Q,W]$.\\

\item The determinant
\begin{equation}\label{eq:intrinsic-det}
 \det\Eul_\delta(t)\in\Z[t,t^{-1}]
\end{equation}
is a well-defined invariant of the dilation class $\delta\in\Gclass[Q,W]$.
\end{enumerate} 
\end{corollary}

\begin{remark}
Lemmas \ref{prop:euler-gauge}, \ref{lem:trivial-cancel} and \ref{lem:right-equivalence} actually imply the stronger statement that the diagonal-monomial conjugacy class of the graded Euler matrix is well-defined for a given dilation class. That is, for any two representatives of the same dilation class, the corresponding graded Euler matrices lie in the same orbit under conjugation by diagonal matrices with monomial entries in the $t$-variable. If vertex relabeling is allowed, the conclusion is modified by using permutation matrices with monomial entries in the $t$-variable, instead of just diagonal matrices.\hfill$\Box$
\end{remark}

%%%%%%%%%%%%%%%%%%%%%%%%%%%%%%%%%%%%%%%%%%%%%%%%%%%%%%%%%%%%%%%%%%%%%
%%%%%%%%%%%%%%%%%%%%%%%%%%%%%%%%%%%%%%%%%%%%%%%%%%%%%%%%%%%%%%%%%%%%%
%%%%%%%%%%%%%%%%%%%%%%%%%%%%%%%%%%%%%%%%%%%%%%%%%%%%%%%%%%%%%%%%%%%%%
\subsection{Graded Euler pairing and graded QP-mutation}\label{sec:euler-mutation}
Consider the involution \(\overline{f(t)}:=f(t^{-1})\) of \(\Lambda=\Z[t,t^{-1}]\), and for a
matrix \(M\in\mbox{GL}(\Lambda)\) set  $M^\dagger=\overline M^{\mathsf T}$. This is needed because the graded Euler pairing is sesquilinear. Indeed, in our convention $[M\langle r\rangle]=t^r[M]$, the graded
Euler pairing satisfies
\[
\chi_t(M\langle r\rangle,N)
=
t^{-r}\chi_t(M,N),
\qquad
\chi_t(M,N\langle r\rangle)
=
t^r\chi_t(M,N).
\]
Thus it is sesquilinear for the involution
$\overline{f(t)}=f(t^{-1})$, as claimed. The matrix for the graded Euler pairing in the simple basis, as in \Cref{def:euler}, behaves as follows under graded QP-mutation:

\begin{proposition}[Effect of graded QP-mutation]\label{prop:mutation-euler}
Let \((Q,W,d)\) be a reduced graded QP, $k\in Q_0$ a mutable vertex and $\mu_k^L(Q,W,d)$ its left graded QP-mutation. Then there is a matrix \(M_k(t)\in\mathrm{GL}_{Q_0}(\Z[t,t^{-1}])\) such that
\begin{equation}\label{eq:mutation-congruence}
 \Eul_{\mu_k^L(Q,W,d)}(t)
 =M_k(t)^\dagger\Eul_{(Q,W,d)}(t)M_k(t)
\end{equation}
In particular, we conclude that
\begin{equation}\label{eq:mutation-coker}
 \coker\Eul_{\mu_k^L(Q,W,d)}(t)
 \cong\coker\Eul_{(Q,W,d)}(t).
\end{equation}
as $\Z[t,t^{-1}]$-modules, and thus also

\begin{equation}\label{eq:mutation-det}
 \det\Eul_{\mu_k^L(Q,W,d)}(t)
 =\det\Eul_{(Q,W,d)}(t).
\end{equation}
\begin{comment}
\begin{equation}\label{eq:det-C}
 \det C_k(t)=\pm t^m
\end{equation}

for some \(m\in\Z\).  Hence
\end{comment}
\end{proposition}

\begin{proof} Let us ease notation by writing $(Q',W',d'):=\mu_k^L(Q,W,d)$ and let \(\Gamma,\Gamma'\) be the corresponding complete differential bigraded Ginzburg algebras. By \cite[Theorem 3.2.(b)]{KellerYang}, QP-mutation induces mutually inverse derived equivalences of $\D_{\cont}(\Gamma)$,
restricting to derived equivalences of the corresponding perfect and finite-dimensional categories. This is achieved by constructing a $(\Gamma,\Gamma')$-bimodule $T$ such that left and right derived tensoring with $T$ induces such derived equivalences, i.e.~they construct a Fourier-Mukai kernel $T$, cf.~\cite[Section 3.4]{KellerYang}. In the case of a graded (homogeneous)
QP-mutation, every generator and differential in their mutation bimodule $T$ is
homogeneous for the internal grading prescribed by
\eqref{eq:mut-composite}--\eqref{eq:mut-outgoing}.  Thus the derived equivalence is
internal-degree preserving.\\

In order to understand the effect on the simples, and the corresponding Euler pairing, we use the mutation triangles for the simples as in
\cite[Remark 3.3]{KellerYang}. Indeed, their derived equivalence for QP-mutation implies that on the graded Grothendieck group
we have
\begin{equation}\label{eq:K0-nonk}
 [F(S_i')]=
 \begin{cases}
[S_i]+p_i(t)[S_k] & \mbox{if }i\ne k\\
-t^m[S_k] & \mbox{if }i=k,
 \end{cases}
\end{equation}
where each $p_i(t)\in\Z[t,t^{-1}]$ is a Laurent polynomial and $m\in\Z$.
\begin{comment}
\(\), while
\begin{equation}\label{eq:K0-k}
 [F(S_k')]=
\end{equation}
for the internal shift fixed by the left-mutation convention.
\end{comment}
The required matrix $M_k(t)\in\Z[t,t^{-1}]$ follows from \Cref{eq:K0-nonk}. Indeed, let \(M_k(t)\)
have these new classes $F(S_i')$ as columns, expressed in the old simple basis. Namely, \(M_k(t)\) is obtained from the identity by adding Laurent multiples of the \(k\)-th basis vector to
other basis vectors and replacing the \(k\)-th basis vector by
\(-t^m\) times itself. Therefore \(\det M_k(t)=-t^m\) and $M_k(t)$ is invertible. Since the derived equivalence in \cite[Theorem 3.2]{KellerYang} preserves the equivariant Euler pairing, by the discussion above, the corresponding matrices satisfy \eqref{eq:mutation-congruence}.\end{proof}

%%%%%%%%%%%%%%%%%%%%%%%%%%%%%%%%%%%%%%%%%%%%%%%%%%%%%%%%%%%%%%%%%%%%%
%%%%%%%%%%%%%%%%%%%%%%%%%%%%%%%%%%%%%%%%%%%%%%%%%%%%%%%%%%%%%%%%%%%%%
%%%%%%%%%%%%%%%%%%%%%%%%%%%%%%%%%%%%%%%%%%%%%%%%%%%%%%%%%%%%%%%%%%%%%
\subsection{A brief summary of perestroikas for the graded Euler form}\label{sec:euler-summary} It might be helpful to summarize the effect on $\Eul_{(Q,W,d)}(t)$ that certain operations on a graded QP have, as follows from \Cref{prop:euler-gauge}, \Cref{lem:trivial-cancel}, \Cref{lem:right-equivalence} and \Cref{prop:mutation-euler}:

\[
\boxed{
\begin{array}{c}
\text{right-equivalence}\\
(Q,W)\rightsquigarrow (Q,W')
\end{array}
\quad\stackrel\Longrightarrow\quad
\quad\Eul_{(Q,W',d)}(t)= \Eul_{(Q,W,d)}(t)
}
\]

\[
\boxed{
\begin{array}{c}
\text{homogeneous stabilization}\\
(Q,W)\rightsquigarrow (Q',W')
\end{array}
\quad\stackrel\Longrightarrow\quad
\quad\Eul_{(Q',W',d)}(t)= \Eul_{(Q,W,d)}(t)
}
\]

\[
\boxed{
\begin{array}{c}
\text{vertex gauging}\\
d\rightsquigarrow d+\delta h
\end{array}
\quad\Longrightarrow\quad
\Eul_{(Q,W,d^h)}(t)= D_h(t)\Eul_{(Q,W,d)}(t)D_h(t)^{-1}
}
\]

\[
\boxed{
\begin{array}{c}
\text{graded QP mutation}\\
(Q,W,d)\rightsquigarrow(Q',W',d')
\end{array}
\quad\Longrightarrow\quad
\Eul_{(Q',W',d')}(t)= M(t)^\dagger \Eul_{(Q,W,d)}(t)M(t)
}
\]

\noindent Therefore, under any of these four operations above, the isomorphism type of the $\Z[t,t^{-1}]$-module $\coker\Eul_{(Q,W,d)}(t)$ remains invariant, see \Cref{rmk:linear_algebra}.

\begin{remark}\label{rmk:linear_algebra} (A commutative algebra reminder) Consider a commutative ring $R$ and a linear map $f:R^n\lr R^m$. Given any $A\in\GL_n(R)$ and $B\in\GL_m(R)$, then the $R$-modules
$$\coker_R(f)\cong \coker_R(A\circ f\circ B)$$
are isomorphic, as $R$-modules. By taking $R=\Z[t,t^{-1}]$ we deduce that both congruence and conjugation each preserves the isomorphism type of the cokernel module $\coker\Eul_{(Q,W,d)}(t)$. For $n=m$, the determinant of $f$ is recovered from the cokernel, up to units of $R$, as the generator of its $0$th Fitting ideal. Since the units of the ring of Laurent polynomials $\Z[t,t^{-1}]$ are $\pm t^{m}$, $m\in\Z$, the determinant $\det\Eul_{(Q,W,d)}(t)$ is recovered from $\coker\Eul_{(Q,W,d)}(t)$ up to a product with $\pm t^{m}$, $m\in\Z$.\hfill$\Box$
\end{remark}

%%%%%%%%%%%%%%%%%%%%%%%%%%%%%%%%%%%%%%%%%%%%%%%%%%%%%%%%%%%%%%%%%%%%%
%%%%%%%%%%%%%%%%%%%%%%%%%%%%%%%%%%%%%%%%%%%%%%%%%%%%%%%%%%%%%%%%%%%%%
%%%%%%%%%%%%%%%%%%%%%%%%%%%%%%%%%%%%%%%%%%%%%%%%%%%%%%%%%%%%%%%%%%%%%
\subsection{Canonically graded quivers}\label{ssec:canonically_gradedQPs} The construction in this section leads to defining the following class of quivers:

\begin{definition}[Canonically graded quivers]\label{def:canonically_graded} A quiver $Q$ is said to be canonically graded if:
\begin{enumerate}
    \item There exists a unique non-degenerate reduced QP $(Q,W)$, up to right-equivalence.

    \item For the unique non-degenerate potential $W$, we have $|\Gclass[Q,W]|=1$.\hfill$\Box$
\end{enumerate}
\end{definition}

\begin{corollary}\label{cor:canonically_graded_invariant}
Let $Q$ be a canonically graded quiver, with its unique non-degenerate reduced QP $(Q,W)$ and dilation $\delta\in \Gclass[Q,W]$. Then:

\begin{enumerate}
    \item The isomorphism type of the $\Z[t,t^{-1}]$-module
    $$\coker \Eul_{(Q,W,\delta)}(t)$$
    is an invariant of the quiver mutation class $[Q]$.\\

    \item The determinant
$$\det\Eul_{(Q,W,\delta)}(t)\in\Z[t,t^{-1}]$$
of the graded Euler pairing is an invariant of the quiver mutation class $[Q]$.
\end{enumerate}
\end{corollary}

\begin{proof}
Let $Q'\in[Q]$ be a quiver mutation equivalent to $Q$. Since QP mutation is invertible and preserves non-degeneracy, $Q$ having a unique non-degenerate potential, up to right equivalence, implies that $Q'$ has a unique non-degenerate potential $W'$, up to right-equivalence. By \Cref{thm:G1-mutation} we obtain that $|\Gclass[Q',W']|=|\Gclass[Q,W]|=1$, and so $\Gclass[Q',W']$ has a unique element $\delta'$. Thus we can name $\det\Eul_{(Q',W',\delta')}(t)$ from $Q'$ intrinsically. \Cref{prop:mutation-euler} then implies the required statements.
%$\det\Eul_{(Q,W,\delta)}(t)=\det\Eul_{(Q',W',\delta')}(t)$, and so the determinant of the graded Euler pairing is an invariant of the quiver mutation class $[Q]$.
\end{proof}

%%%%%%%%%%%%%%%%%%%%%%%%%%%%%%%%%%%%%%%%%%%%%%%%%%%%%%%%%%%%%%%%%%%%%
%%%%%%%%%%%%%%%%%%%%%%%%%%%%%%%%%%%%%%%%%%%%%%%%%%%%%%%%%%%%%%%%%%%%%
%%%%%%%%%%%%%%%%%%%%%%%%%%%%%%%%%%%%%%%%%%%%%%%%%%%%%%%%%%%%%%%%%%%%%
\subsection{The example of tree quivers}\label{ssec:tree_quivers} Let $(Q,W)$ be a reduced QP and $\Gclass[Q,W]$ its set of dilation classes. By the results of \Cref{sec:euler-dilation} and \Cref{sec:euler-mutation}, the set of polynomials
$$\Pclass(Q,W):=\{\det\Eul_\delta(t)\in\Z[t,t^{-1}]:\mbox{ }\delta\in\Gclass[Q,W]\}$$
is an invariant of the QP-mutation class of $(Q,W)$. In particular:\\

\begin{enumerate}
    \item If $|\Gclass[Q,W]|=1$, i.e.~if there exists a unique dilation class, then $\Pclass(Q,W)$ is a unique Laurent polynomial. Thus we obtain in this case a Laurent polynomial as an invariant of a QP-mutation class.\\

    \item If $|\Gclass[Q,W]|=1$ and $W$ is the unique non-degenerate potential for $Q$, up to right equivalence, as in \Cref{def:canonically_graded}, then we obtain in this case a Laurent polynomial as an invariant of a quiver mutation class, cf.~\Cref{cor:canonically_graded_invariant}. Namely, for each quiver in the mutation class we can select the right-equivalence class of the unique non-degenerate potential and consider its unique dilation. From there $\Pclass(Q,W)$ provides the Laurent polynomial invariant of the quiver mutation class.\\
\end{enumerate}

\noindent There is a first example where $(2)$ above holds: for any quiver $Q$ whose underlying graph is a tree, cf.~\cite[Def.~1.6]{FominNeville}. That is, for any quiver obtained by considering a tree and simply orienting its edges. Though not logically needed for our main result, we briefly study this case of tree quivers for completeness. For a tree quiver, since there are no cycles, any potential is automatically zero, and the arrow-gradings are unique up to vertex gauge. More precisely:

\begin{lemma}[$\exists!$ reduced graded QP for tree quivers]\label{lem:unique-tree-qp}
Let $Q$ be a tree quiver. Then:
\begin{enumerate}[label=\textup{(\roman*)}]
\item There exists a unique potential $W=0$ on $Q$ and the corresponding QP is rigid, and hence non-degenerate.\\

\item Every integral arrow grading $d\colon Q_1\to\mathbb Z$ is compatible with the requirement that $W$ is homogeneous of degree $1$, and all such gradings are equivalent under vertex gauge. Thus the unique gauge class is represented by the zero grading $d(a)=0$ for all arrows $a\in Q_1$. 
\end{enumerate}
\end{lemma}

\begin{proof}
For $(i)$, since the underlying unoriented graph is a tree, $Q$ has no unoriented cycles and thus $W=0$ is the only possible potential. As there are therefore no nonzero cyclic deformations, the QP $(Q,0)$ is rigid and, by \cite[Corollary~8.2]{DWZ}, non-degenerate.\\

\noindent For $(ii)$, since the zero potential is homogeneous of every degree, every arrow-grading $d\colon Q_1\to\mathbb Z$ is compatible. Recall that a vertex gauge is a function $h\colon Q_0\to\mathbb Z$ acting on an arrow-grading $d$ by
\[
 d^{h}(a)=d(a)+h(t(a))-h(s(a)).
\]
Thus the quotient of the $\Z$-span of the arrows by vertex gauging is canonically $H^1(|Q|;\mathbb Z)$. Since $|Q|$ is a tree, $H^1(|Q|;\mathbb Z)=0$.  Thus all arrow gradings are gauge-equivalent, and the zero grading is a representative.  Homogeneous trivial QP summands disappear under graded QP-reduction, and so they do not alter the reduced stable class.
\end{proof}

\begin{remark}
The conclusion in \Cref{lem:unique-tree-qp}.(ii) uses that $Q$ is a tree quiver, not just an acyclic quiver. For instance, the quiver $Q$ with three vertices $\{1,2,3\}$ and arrows $1\to2$, $2\to3$ and $1\to3$ is acyclic but $H^1(|Q|;\Z)\cong\Z$ does not vanish. Thus the argument fails and, in fact, there is not a unique dilation class in this case. Similarly, \Cref{lem:matching} below only holds for tree quivers, but not for general acyclic quivers.\hfill$\Box$
\end{remark}

\noindent By \Cref{prop:euler-formula}, and using the zero potential and grading as in \Cref{lem:unique-tree-qp}, the graded Euler pairing $\cE_Q(t):= \cE_{(Q,0,0)}(t)$ in terms of the simple modules of the differential bigraded Ginzburg algebra reads combinatorially as:

\begin{equation}\label{eq:graded-euler-entry}
 \boxed{
 \cE_Q(t)_{ij}
 =(1-t)\delta_{ij}
 -\#\{a:j\to i\}
 +t\,\#\{a:i\to j\}.}
\end{equation}
Equivalently, every arrow $i\to j$ contributes $t$ in position $(i,j)$ and
$-1$ in position $(j,i)$, while every diagonal entry is $1-t$. From \eqref{eq:graded-euler-entry} there is a matching expansion for the determinant in certain cases, as follows:

\begin{lemma}[Matching expansion]\label{lem:matching}
Let $T$ be a tree quiver with $N$ vertices and no multiple arrows, and let $m_k(T)$ be the number of matchings of $T$ with exactly $k$ edges.  Then
\begin{equation}\label{eq:matching-formula}
 \boxed{
 \det\cE_T(t)=\sum_{k\ge 0}m_k(T)t^k(1-t)^{N-2k}.
 }
\end{equation}
In particular, the determinant is independent of the orientation of the tree.
\end{lemma}

\begin{proof}
First, expand the determinant as a sum over permutations. If a nonzero permutation term contained a cycle of length at least three, the corresponding nonzero off-diagonal matrix entries would exhibit an unoriented cycle in $T$. This is impossible because $T$ is a tree and thus every contributing permutation is a product of disjoint transpositions along edges and fixed points. In consequence, its transposed edges form a matching. For each transposed edge $\{i,j\}$, the permutation sign contributes $-1$, whereas the two matrix entries are $t$ and $-1$ (in one order or the other), and hence $\cE_{ij}(t)\cE_{ji}(t)=t(-1)=-t$.  Their combined contribution is therefore $t$. Since each fixed point contributes $1-t$, a matching of size $k$ thus contributes $t^k(1-t)^{N-2k}$. By summing over all matchings we obtain \eqref{eq:matching-formula}.
\end{proof}

%%%%%%%%%%%%%%%%%%%%%%%%%%%%%%%%%%%%%%%%%%%%%%%%%%%%%%%%%%%%%%%%%%%%%%%%%
%%%%%%%%%%%%%%%%%%%%%%%%%%%%%%%%%%%%%%%%%%%%%%%%%%%%%%%%%%%%%%%%%%%%%%%%%
\subsubsection{Relation to \cite{FominNeville}} In the article \cite{FominNeville}, the authors use cyclically ordered quivers to construct polynomials that are invariant under proper mutation of a cyclically ordered quiver. The case of a tree quiver $Q$ is discussed in detail in \cite[Section 10]{FominNeville}: by \cite[Prop.~2.6\&Cor.~4.11]{FominNeville} their polynomial $\FNDelta_Q(t)$ is independent of the choice of cyclic ordering. Let us now compare $\det\cE_Q(t)$, as introduced in \Cref{def:euler}, to $\FNDelta_Q(t)$, from \cite{FominNeville}, as follows.\\

\noindent Given a tree quiver $Q$, let us consider the matrix $B_Q$ with entries
\[
 b_{ij}:=\#\{i\to j\}-\#\{j\to i\},
\]
and let $U_Q$ be the unipotent companion, cf.~\cite[Def.~4.2]{FominNeville}, given by the unique upper-triangular matrix with diagonal entries $1$ satisfying
\begin{equation}\label{eq:FN-companion}
 -B_Q=U_Q-U_Q^{\mathsf T}.
\end{equation}
By construction, see \cite[Def.~8.1]{FominNeville}, we have $\FNDelta_Q(t)=\det(tU_Q-U_Q^{\mathsf T})$. The comparison then reads:

\begin{lemma}[Comparison to \cite{FominNeville}]\label{lem:FN-comparison}
Let $Q$ be a tree quiver on $N$ vertices, and choose an ordering of its vertices so that every arrow points from a smaller vertex to a larger one. Then
\begin{equation}\label{eq:E-U-trees}
 \cE_{(Q,0,0)}(t)=U_Q^{\mathsf T}-tU_Q.
\end{equation}
In consequence, we obtain the equality of polynomials
\begin{equation}\label{eq:sign-comparison}
 \boxed{\det\cE_{(Q,0,0)}(t)=(-1)^N\FNDelta_Q(t).}
\end{equation}
\end{lemma}

\begin{proof}
For two vertices $i,j\in Q_0$ with $i<j$, the ordering implies that there are no arrows $j\to i$.  By \eqref{eq:FN-companion},
\[
 (U_Q)_{ij}=-b_{ij}=-\#\{i\to j\},
 \qquad (U_Q)_{ji}=0.
\]
Thus the diagonal entries of $U_Q^{\mathsf T}-tU_Q$ are $1-t$.  For every arrow $i\to j$, the $(i,j)$-entry is $-t(U_Q)_{ij}=t$ whereas the $(j,i)$-entry is
$(U_Q^{\mathsf T})_{ji}=(U_Q)_{ij}=-1$. This yields precisely \eqref{eq:graded-euler-entry} and it implies \eqref{eq:E-U-trees}. For the determinants, note that \eqref{eq:E-U-trees} gives
\[
 \FNDelta_Q(t)
 =\det(tU_Q-U_Q^{\mathsf T})
 =\det(-\cE_{(Q,0,0)}(t))
 =(-1)^N\det\cE_{(Q,0,0)}(t),
\]
which is equivalent to \eqref{eq:sign-comparison}.
\end{proof}

\begin{remark}
Note that all orientations of a fixed tree quiver are mutation-equivalent. In addition, \cite[Theorem 1.1]{Neville} states that mutation-acyclic quivers are totally proper, thus promoting the polynomial $\FNDelta_Q(t)$ to an invariant of any quiver mutation class containing a tree. From this viewpoint, Lemma~\ref{lem:FN-comparison} identifies that invariant with a normalization of the determinant of the graded Euler pairing in the category of finite-dimensional continuous differential bigraded modules over the differential bigraded Ginzburg algebra of $(Q,0)$.\hfill$\Box$
\end{remark}

\noindent We conclude this subsection with a few computations to illustrate the results above.

%%%%%%%%%%%%%%%%%%%%%%%%%%%%%%%%%%%%%%%%%%%%%%%%%%%%%%%%%%%%%%%%%%%%%%%%%
%%%%%%%%%%%%%%%%%%%%%%%%%%%%%%%%%%%%%%%%%%%%%%%%%%%%%%%%%%%%%%%%%%%%%%%%%

\subsubsection{The linear tree \texorpdfstring{$A_{2g}$}{A2g}}

Let $g\in\N$ and consider the linearly oriented tree on $2g$ vertices:
\[
\begin{tikzpicture}[baseline=(current bounding box.center),node distance=12mm]
  \node[qvertex] (v1) {$1$};
  \node[qvertex,right=of v1] (v2) {$2$};
  \node[qvertex,right=of v2] (v3) {$3$};
  \node[right=10mm of v3] (dots) {$\cdots$};
  \node[qvertex,right=10mm of dots] (vn1) {$2g-1$};
  \node[qvertex,right=of vn1] (vn) {$2g$};
  \draw[qarrow] (v1)--(v2);
  \draw[qarrow] (v2)--(v3);
  \draw[qarrow] (v3)--(dots);
  \draw[qarrow] (dots)--(vn1);
  \draw[qarrow] (vn1)--(vn);
\end{tikzpicture}
\]
By \eqref{eq:euler-formula} or \eqref{eq:graded-euler-entry}, the graded Euler matrix from \Cref{def:euler} is the tridiagonal matrix
\begin{equation}\label{eq:A-matrix}
 \cE_{A_{2g}}(t)=
 \begin{pmatrix}
 1-t&t&0&\cdots&0\\
 -1&1-t&t&\ddots&\vdots\\
 0&-1&1-t&\ddots&0\\
 \vdots&\ddots&\ddots&\ddots&t\\
 0&\cdots&0&-1&1-t
 \end{pmatrix}.
\end{equation}
Let $D_m(t)$ denote the determinant of the analogous $m\times m$ matrix.  The product of the two off-diagonal entries adjacent to the last diagonal entry is $t(-1)=-t$, so expansion along the last row or column gives
\[
 D_m(t)=(1-t)D_{m-1}(t)+tD_{m-2}(t),
 \qquad D_0(t)=1,\quad D_1(t)=1-t.
\]
Resolving this recursion, e.g.~via induction, yields
\[
 D_m(t)=\sum_{j=0}^{m}(-t)^j.
\]
and thus for $m=2g$ we obtain
\begin{equation}\label{eq:A-det}
 \boxed{
 \det\cE_{A_{2g}}(t)
 =1-t+t^2-\cdots+t^{2g}
 =\frac{1+t^{2g+1}}{1+t}.
 }
\end{equation}

Note that the irreducible plane curve singularity $x^2+y^{2g+1}=0$ has the $(2,2g+1)$-torus knot as its knot, obtained by intersecting with the 3-sphere boundary of a small enough Milnor ball. The Alexander polynomial $\Delta_{T(2,2g+1)}(t)$ of such a torus knot is
\[
 \Delta_{T(2,2g+1)}(t)
 =\frac{(1-t)(1-t^{4g+2})}{(1-t^2)(1-t^{2g+1})}
 =\frac{1+t^{2g+1}}{1+t}.
\]
Thus
\[
 \boxed{
 \det\cE_{A_{2g}}(t)=\Delta_{T(2,2g+1)}(t).
 }
\]
\noindent In subsequent sections, we will be generalizing such equality for any irreducible plane curve singularity that admits a real Morsification with a malleable divide, where the quiver is no longer necessarily a tree quiver.

\subsubsection{\texorpdfstring{$E_6$}{E6} tree}

Choose the following orientation and labeling:
\[
\begin{tikzpicture}[baseline=(current bounding box.center),node distance=11mm]
  \node[qvertex] (v1) {$1$};
  \node[qvertex,right=of v1] (v2) {$2$};
  \node[qvertex,right=of v2] (v3) {$3$};
  \node[qvertex,right=of v3] (v4) {$4$};
  \node[qvertex,right=of v4] (v5) {$5$};
  \node[qvertex,below=10mm of v3] (v6) {$6$};
  \draw[qarrow] (v1)--(v2);
  \draw[qarrow] (v2)--(v3);
  \draw[qarrow] (v3)--(v4);
  \draw[qarrow] (v4)--(v5);
  \draw[qarrow] (v3)--(v6);
\end{tikzpicture}
\]
By \eqref{eq:graded-euler-entry} the graded Euler matrix is
\begin{equation}\label{eq:E6-matrix}
 \cE_{E_6}(t)=
 \begin{pmatrix}
 1-t&t&0&0&0&0\\
 -1&1-t&t&0&0&0\\
 0&-1&1-t&t&0&t\\
 0&0&-1&1-t&t&0\\
 0&0&0&-1&1-t&0\\
 0&0&-1&0&0&1-t
 \end{pmatrix}.
\end{equation}
\noindent For the purpose of illustration, let us compute its determinant via matchings as in \Cref{lem:matching}. In this case the non-zero matching numbers are $(m_0,m_1,m_2,m_3)=(1,5,5,1)$ and \Cref{lem:matching} gives
\begin{align*}
 \det\cE_{E_6}(t)
 &=(1-t)^6+5t(1-t)^4+5t^2(1-t)^2+t^3\\
 &=t^6-t^5+t^3-t+1.
\end{align*}
Therefore
\begin{equation}\label{eq:E6-det}
 \boxed{
 \det\cE_{E_6}(t)=t^6-t^5+t^3-t+1
 =(t^2-t+1)(t^4-t^2+1).
 }
\end{equation}
As in the $A_{2g}$ case above, the irreducible plane curve singularity $x^3+y^{4}=0$ has the $(3,4)$-torus knot as the algebraic knot of the singularity. Its Alexander polynomial $\Delta_{T(3,4)}(t)$ is
\[
 \Delta_{T(3,4)}(t)
 =\frac{(1-t)(1-t^{12})}{(1-t^3)(1-t^4)}
 =t^6-t^5+t^3-t+1.
\]
Hence
\[
 \boxed{\det\cE_{E_6}(t)=\Delta_{T(3,4)}(t).}
\]

\subsubsection{A few more instances} Consider the tree quivers in Figures \ref{fig:E8}, \ref{fig:tree1} and \ref{fig:tree2}. The determinants $\det\cE_T(t)$ of their graded Euler pairings $\cE_T(t)$ are computed in \Cref{table:tree_quivers}.

\begin{center}
\small
\renewcommand{\arraystretch}{1.35}
\begin{table}[htbp]
    \centering
\caption{The determinants $\det\cE_T(t)$ of the graded Euler pairings for a few tree quivers.}
\label{table:tree_quivers}
\begin{tabularx}{\textwidth}{>{\raggedright\arraybackslash}p{1.55cm}
>{\raggedright\arraybackslash}X
>{\raggedright\arraybackslash}p{3.65cm}
>{\raggedright\arraybackslash}p{3.05cm}}
\toprule
Tree & $\det\cE_T(t)$ & Alexander poly.~ of & Algebraic?\\
\midrule
$A_{2g}$
& $1-t+t^2-\cdots+t^{2g}$
& $T(2,2g+1)$
& Yes: $x^2=y^{2g+1}$\\
$E_6$
& $t^6-t^5+t^3-t+1$
& $T(3,4)$
& Yes: $x^3=y^4$\\
$E_8$
& $t^8-t^7+t^5-t^4+t^3-t+1$
& $T(3,5)$
& Yes: $x^3=y^5$\\
$T_{(2,2,3)}$
& $t^8-t^7+2t^5-3t^4+2t^3-t+1$
& Slalom knot $K_{(2,2,3)}$
& No\\
$T_{(4,4,1)}$
& $t^{10}-t^9+t^7-2t^6+3t^5-2t^4+t^3-t+1$
& Slalom knot $K_{(4,4,1)}$
& No\\
\bottomrule
\end{tabularx}
\end{table}
\end{center}

\begin{figure}
\centering
\[
\begin{tikzpicture}[baseline=(current bounding box.center),node distance=9.5mm]
  \node[qvertex] (v1) {$1$};
  \node[qvertex,right=of v1] (v2) {$2$};
  \node[qvertex,right=of v2] (v3) {$3$};
  \node[qvertex,right=of v3] (v4) {$4$};
  \node[qvertex,right=of v4] (v5) {$5$};
  \node[qvertex,right=of v5] (v6) {$6$};
  \node[qvertex,right=of v6] (v7) {$7$};
  \node[qvertex,below=10mm of v3] (v8) {$8$};
  \draw[qarrow] (v1)--(v2);
  \draw[qarrow] (v2)--(v3);
  \draw[qarrow] (v3)--(v4);
  \draw[qarrow] (v4)--(v5);
  \draw[qarrow] (v5)--(v6);
  \draw[qarrow] (v6)--(v7);
  \draw[qarrow] (v3)--(v8);
\end{tikzpicture}
\]
\caption{The $E_8$ tree with a choice of orientation. The associated $\det\Eul_Q(t)$ gives the Alexander polynomial of the $(3,5)$-torus knot, associated to the plane curve singularity $x^3+y^5=0$.}
\label{fig:E8}
\end{figure}

\begin{figure}
\centering
\[
\begin{tikzpicture}[baseline=(current bounding box.center),node distance=10.5mm]
  \node[qvertex] (v1) {$1$};
  \node[qvertex,right=of v1] (v2) {$2$};
  \node[qvertex,right=of v2] (v3) {$3$};
  \node[qvertex,right=of v3] (v4) {$4$};
  \node[qvertex,right=of v4] (v5) {$5$};
  \node[qvertex,below=9mm of v3] (v6) {$6$};
  \node[qvertex,below=9mm of v6] (v7) {$7$};
  \node[qvertex,below=9mm of v7] (v8) {$8$};
  \draw[qarrow] (v1)--(v2);
  \draw[qarrow] (v2)--(v3);
  \draw[qarrow] (v3)--(v4);
  \draw[qarrow] (v4)--(v5);
  \draw[qarrow] (v3)--(v6);
  \draw[qarrow] (v6)--(v7);
  \draw[qarrow] (v7)--(v8);
\end{tikzpicture}
\]
\caption{A tree quiver $T_{(2,2,3)}$, with a central trivalent vertex and three arms of length $(2,2,3)$. In this case $\det\cE_{T_{(2,2,3)}}(t)$ is not the Alexander polynomial of any algebraic knot, by the semi-group gap obstruction, cf.~\cite[Section~3.1]{BorodzikHom} or \cite{CDGZ}. It is the Alexander polynomial of the slalom knot $K_{(2,2,3)}$, cf.~\cite{ACampo}, which is a fibered hyperbolic knot.}
\label{fig:tree1}
\end{figure}

\begin{figure}
\centering
\[
\begin{tikzpicture}[baseline=(current bounding box.center),node distance=8.2mm]
  \node[qvertex] (v1) {$1$};
  \node[qvertex,right=of v1] (v2) {$2$};
  \node[qvertex,right=of v2] (v3) {$3$};
  \node[qvertex,right=of v3] (v4) {$4$};
  \node[qvertex,right=of v4] (v5) {$5$};
  \node[qvertex,right=of v5] (v6) {$6$};
  \node[qvertex,right=of v6] (v7) {$7$};
  \node[qvertex,right=of v7] (v8) {$8$};
  \node[qvertex,right=of v8] (v9) {$9$};
  \node[qvertex,below=10mm of v5] (v10) {$10$};
  \draw[qarrow] (v1)--(v2);
  \draw[qarrow] (v2)--(v3);
  \draw[qarrow] (v3)--(v4);
  \draw[qarrow] (v4)--(v5);
  \draw[qarrow] (v5)--(v6);
  \draw[qarrow] (v6)--(v7);
  \draw[qarrow] (v7)--(v8);
  \draw[qarrow] (v8)--(v9);
  \draw[qarrow] (v5)--(v10);
\end{tikzpicture}
\]
\caption{A tree quiver $T_{(4,4,1)}$, with a central trivalent vertex and three arms of length $(4,4,1)$. In this case $\det\cE_{T_{(4,4,1)}}(t)$ is also not the Alexander polynomial of any algebraic knot, by the same semi-group obstruction. It is the Alexander polynomial of the slalom knot $K_{(4,4,1)}$, a genus-five fibered hyperbolic knot.}
\label{fig:tree2}
\end{figure}